\documentclass[11pt,a4paper]{article}
\usepackage[margin=26mm]{geometry}
\usepackage[T1]{fontenc}
\usepackage{lmodern,microtype}
\usepackage{amsmath,amssymb,amsthm,mathtools,booktabs}
\usepackage[hidelinks]{hyperref}
\hypersetup{pdftitle={Joint Minorants for Nine Points},
 pdfauthor={RH-Weil research project}}
\newtheorem{theorem}{Theorem}[section]
\newtheorem{lemma}[theorem]{Lemma}
\newtheorem{proposition}[theorem]{Proposition}
\newtheorem{inputclaim}[theorem]{Computational input}
\theoremstyle{remark}
\newtheorem{remark}[theorem]{Remark}
\DeclareMathOperator{\tr}{tr}
\DeclareMathOperator{\sinc}{sinc}
\newcommand{\R}{\mathbb R}
\newcommand{\I}{[-1/2,1/2]}
\title{Joint Minorants for Nine Points\\
and Simple Critical Zeros of the Riemann Zeta Function}
\author{RH--Weil research project}
\date{October 8, 2026}
\begin{document}
\maketitle
\begin{abstract}
We improve a fixed eight-point certificate for a perturbed
Montgomery--Taylor window by jointly treating two overlapping frames.
The two frames share their actual distance-square variables and all
common gap and span constraints. From the existing continuous certificate
we capture 70 low-slack cells, restore their reflections, and reduce the
joint problem to 289 labelled intersections. Rational epigraph duals,
checked with exact finite-box residual corrections, pay the local reward
$805103/10^8$ on every intersection. Together with the existing
lossless separation and all-zero spectral assembly this gives
\[
 \liminf_{T\to\infty}\frac{N_0^s(T)}{N(T)}
 \ge\frac{66812491}{99194897}=0.6735476624367078\ldots.
\]
The denominator counts all nontrivial zeros with multiplicity; the
numerator counts simple zeros on the critical line. Off-line zeros are
retained. The analytic argument assumes no zero-free strip or RH.
The finite computational input retains the disclosed PC8CL
continuous-audit and compiled-integer trust scope. The new dual verifier
uses only integer and rational arithmetic; numerical linear programming
merely locates witnesses. This is not a complete Lean kernel proof.
\end{abstract}

\section{Statement, provenance and comparison}

Let $N(T)$ count the nontrivial zeros $\rho$ of $\zeta$, with
$0<\Im\rho\le T$, including multiplicity. Let $N_0^s(T)$ count the
simple zeros among them with $\Re\rho=1/2$.

\begin{theorem}\label{thm:main}
Using the continuous PC8CL certificate stated in
Input~\ref{in:pc8} and the exact joint certificate established below,
one has, unconditionally,
\begin{equation}\label{eq:main}
 \liminf_{T\to\infty}\frac{N_0^s(T)}{N(T)}
 \ge \frac{66812491}{99194897}.
\end{equation}
No zero-free strip, Riemann hypothesis or higher correlation
conjecture is assumed.
\end{theorem}

The window and eight-point data are prior inputs from the fixed
public AM construction \cite{PC8}. Their lossless assembly
\cite{Lossless} gave $66812491/99194997$. The present improvement is
in the local certificate: two overlapping frames share the actual kernel
squares and all their continuous linear span constraints. We adapt the
separation mechanism used by Knausg{\aa}rd \cite[Section~5]{Knaus},
without consuming that paper's headline native computation axiom.
The new rational bound exceeds the previous repository bound by
\[
 \frac{6681249100}{99194897\cdot99194997}>0,
\]
approximately $0.0000679014$ percentage points. Its comparison with the
fixed Knausg{\aa}rd v1 benchmark is
\begin{equation}\label{eq:comparison}
 \frac{66812491}{99194897}-\frac{1669159}{2478195}
 =\frac{2326052122}{245824297770915}>0.
\end{equation}
This is a fixed-source comparison, not an exhaustive global priority
claim. A zero-counting proportion is not a measure of progress toward a
proof of RH, and the separate zero-free-region bounds are unchanged.

\section{Window and admitted local certificate}

Put $\beta=1/\sqrt2$ and $Z_0=\sqrt2\sin\beta$.
On $I=\I$ define
\begin{equation}\label{eq:window}
 f(u)=\frac{\cos(\sqrt2u)+\sum_{j=1}^{12}c_j\cos(2\pi ju)}{Z_0},
 \qquad K(t)=\int_I f(u)e^{-2\pi itu}\,du.
\end{equation}
Extend $f$ by zero outside $I$. The integer numerators $10^9c_j$ are
\[
\begin{split}
(&12310798,-15041681,3867664,6489926,-992327,5580097,\\
 &-6846472,3781297,-5670353,3355089,-450523,-218483).
\end{split}
\]
Thus $S=\sum_j|c_j|=6460471/10^8<13/200$.
Since $\cos(\sqrt2u)\ge1-u^2\ge3/4$ on $I$, the numerator
is positive. The integer cosine modes integrate to zero, so
$f$ is an even probability density. Also
\begin{equation}\label{eq:normalization}
 Z_0=\int_I\cos(\sqrt2u)\,du\ge 11/12.
\end{equation}

For eight ordered points $y_0,\ldots,y_7$, write
$g_r=y_{r+1}-y_r\ge0$, $0\le r\le6$. Set
\[
 c_0=\frac{805003}{10^8},\quad c=\frac{805103}{10^8},\qquad B=\sum_rb_r=\frac{404350}{10^8},
\]
\[
 (10^8b_r)_{r=0}^6=(28898,57272,75526,80958,75526,57272,28898).
\]
All pair weights $a_{ij}$ are nonnegative. Table~\ref{tab:weights}
gives their integer numerators.

\begin{table}[ht]
\centering\small
\begin{tabular}{c*{7}{r}}
\toprule
$i\backslash j$&1&2&3&4&5&6&7\\
\midrule
0&13630830&0&19565904&20317441&45030671&100000000&200000000\\
1&&29997996&30621878&47766489&79682558&109938656&100000000\\
2&&&37356140&69378121&65335210&79682558&45030671\\
3&&&&38030065&69378121&47766489&20317441\\
4&&&&&37356140&30621878&19565904\\
5&&&&&&29997996&0\\
6&&&&&&&13630830\\
\bottomrule
\end{tabular}
\caption{The weights $10^8a_{ij}$ of the fixed AM certificate.}
\label{tab:weights}
\end{table}

\begin{inputclaim}[Fixed PC8CL certificate]\label{in:pc8}
For every $(g_0,\ldots,g_6)\in[0,\infty)^7$,
\begin{equation}\label{eq:local}
 \sum_{r=0}^6b_rg_r+
 \sum_{0\le i<j\le7}a_{ij}
 K\left(\sum_{r=i}^{j-1}g_r\right)^2\ge c_0.
\end{equation}
\end{inputclaim}

This is a continuous real-gap assertion, not a grid sampling claim.
Its provenance, finite replay and continuous soundness evidence are
specified in Section~\ref{sec:trust}. We use the complete admitted
certificate, rather than the upstream final zeta theorem.
Direct integer summation of Table~\ref{tab:weights} gives
\begin{equation}\label{eq:capacity}
\begin{split}
10^8\left(\sum_{i=0}^{7-s}a_{i,i+s}\right)_{s=1}^7
={}&(199999997,199999998,199999996,199999998,\\
 &199999998,200000000,200000000).
\end{split}
\end{equation}
In particular each span-length budget is at most two.

\section{The joint continuous nine-point certificate}\label{sec:joint}

Write $F_8(h)$ for the left side of \eqref{eq:local} and put
$\theta=4/5$, $\delta=10^{-6}$, so $c=c_0+\delta$.
For $g=(g_0,\ldots,g_7)\in[\theta,\infty)^8$, let
$h^-=(g_0,\ldots,g_6)$ and $h^+=(g_1,\ldots,g_7)$. Define
\begin{equation}\label{eq:nine}
 F_9(g)=\tfrac12\{F_8(h^-)+F_8(h^+)\}
                +2K(g_0+\cdots+g_7)^2.
\end{equation}
The new pair weights are
\[
 a'_{ij}=\tfrac12a_{ij}{\bf1}_{j\le7}
        +\tfrac12a_{i-1,j-1}{\bf1}_{i\ge1}
        +2{\bf1}_{(i,j)=(0,8)},
\]
where out-of-range old weights are zero. The gap fees are
\[
 (10^8b'_r)_{r=0}^7=(14449,43085,66399,78242,
                         78242,66399,43085,14449).
\]
All weights and fees are nonnegative; $\sum b'_r=B$.
Each index-span budget is at most two, including span eight.

\subsection{A complete low-slack outer cover}

Put $\mathcal L=\{h\in[\theta,\infty)^7:F_8(h)<c_0+2\delta\}$.
If either frame is outside $\mathcal L$, Input~\ref{in:pc8} and
\eqref{eq:nine} already give $F_9\ge c$.
We explain why the captured cells cover all of $\mathcal L$.
The old small-gap domain $h_i\le1/2$ is absent. In any old clear
COV interval, a single adjacent-pair term already pays $c_0$;
the pressure adds at least $\theta B=8087/2500000>2\delta$.
In any old large-gap domain, one fee pays $c_0$ and the remaining
fees add at least $\theta(B-\max b_i)=5053/1953125>2\delta$.
Thus only the bad COV intervals and their 32 root trees remain.

In the half-space $h_6\ge h_0$, the old direction leaves
$u_6<l_0$ cannot contain the point. Farkas empty-cell leaves do
not contain it either. Every remaining successful leaf supplies an
interval minorant. Keep the original reward, return value and packed
bit cursor; copy only its final reward comparison with the new target
$805203/10^8=c_0+2\delta$. A passing stronger leaf excludes the
point from $\mathcal L$; a nonpassing stronger leaf with a passing
old comparison is captured. The actual unchanged 32 root consumption
positions and seven outer COV checks all pass. This captures exactly
70 cells. Their full reflections, using
$h_r\mapsto h_{6-r}$ and $(i,j)\mapsto(7-j,7-i)$ on spans, restore
the other half-space. There are 140 source labels and 132 distinct
geometric cells. These are outer cells, not subsets on which every
point necessarily has low slack.

Each cell retains seven gap and 21 long-span closed bounds, in units
$SC=32768$. A left/right intersection retains all eight gaps and
27 distinct long spans. Prefix coordinates turn these into exact
closed difference constraints. In units $5SC$ the lower gap bound
$\theta$ is exactly the integer 131072. Negative cycles exclude empty
intersections; shortest-path potentials verify feasibility otherwise.
The 19600 labelled pairs reduce to 422 after common gap bounds and
to 289 after all common spans and the complete closure. Labels are
retained even when their geometric cells coincide, since their
minorants can differ. There is no integer rounding of a real gap.

The capture and continuous semantics reuse the admitted fixed PC8
input. The unchanged old reward comparisons, root cursors and atom
validity checks are essential. A failed stronger comparison alone
would not prove a counterexample, and the list of cells alone would
not prove its completeness.

\subsection{Valid individual tangent cuts}

Let $w(t)=K(t)^2$. For an enabled point $p\in[l,u]$, the old
continuous tangent assertion gives real $v,d$ with
\[
 w(x)\ge v+d(x-p)\quad(l\le x\le u),\qquad
 v\ge v_0,\quad d_-\le d\le d_+.
\]
The enabled-point condition and interval-valid tangent are part of
the admitted value, derivative and convexity or extrapolation checks.
A derivative interval at a point alone would not establish this assertion.
Its packed integer data provide $v_0=V/10^{10}$,
$d_-=(P-2\cdot10^9)/10^9$ and $d_+=(2\cdot10^9-M)/10^9$.
In particular $d_-\le d_+$.

\begin{lemma}\label{lem:cuts}
The following two affine functions lie below $w$ on $[l,u]$:
\[
 \ell_p(x)=v_0+d_+(l-p)+d_-(x-l),\qquad
 r_p(x)=v_0+d_-(u-p)+d_+(x-u).
\]
Any constant bound $\lambda\le\inf_{[l,u]}w$ can also be used.
\end{lemma}
\begin{proof}
Subtracting the first line from $v_0+d(x-p)$ gives
$(d-d_-)(x-l)+(d_+-d)(p-l)\ge0$.
Subtracting the second gives
$(d_+-d)(u-x)+(d-d_-)(u-p)\ge0$.
Both conclusions use $p\in[l,u]$. Only points with positive weight
in the old certified minorant are enabled; a zero-weight point does
not automatically supply a tangent.
\end{proof}

For each captured term the true endpoints are recovered from the
original gap and span constraints. In a zero-width query the value
pyramid may use an enlarged one-grid-cell interval; its constant bound
remains valid, but the true endpoint is used for the tangent cuts.
Every old encoded derivative error and value rounding is retained.
All physical distances shared by the two frames use the same variable
$z_{ij}=w(g_i+\cdots+g_{j-1})$. The two frames have 33 old nonzero
physical pair terms; the added span-eight term makes 34 variables.
The pointwise positivity of $f$ gives $0\le z_{ij}\le1$.
Place all constant and enabled individual cuts from both frames on
that same variable, rather than retaining only the original mixture
or an independently minimized frame constant.
Reflection transports each term by its physical span; it does not
reverse the list of 26 encoded terms blindly.

\subsection{Exact finite-box dual verification}

For each of the 289 intersections form a vector $x$ of eight gaps
and 34 distance-square variables. Write all cell and tangent cuts as
$Ax\le b$, and let $\ell\le x\le u$ be the finite variable box:
closed difference-constraint projections for the gaps, and $[0,1]$
for the squares. The objective vector $q$ represents $2\cdot10^8F_9$;
its gap coefficients are the integer numerators of twice $b'$, and
its pair coefficients are the sum of the old integer weights, with
$4\cdot10^8$ on $(0,8)$. Every actual assignment satisfies all these
constraints and gives the exact objective.

\begin{lemma}[Residual-corrected dual]\label{lem:dual}
For any rational $\lambda\ge0$, put $r=q+A^t\lambda$. Then
\[
 q^tx\ge-\lambda^tb+
 \sum_i r_i\begin{cases}\ell_i,&r_i\ge0,\\u_i,&r_i<0.\end{cases}
 \qquad(Ax\le b,\ \ell\le x\le u).
\]
\end{lemma}
\begin{proof}
Use $q^tx=r^tx-\lambda^tAx\ge r^tx-\lambda^tb$,
and minimize each residual coefficient on its complete box.
Exact stationarity $r=0$ and floating-point feasibility are unnecessary.
\end{proof}

Numerical linear programming only proposes $\lambda$. The certificate
stores its nonnegative rational entries; the verifier reconstructs every
matrix, objective and box from the captured integers and consumes the
entries by Lemma~\ref{lem:dual}. Across all 289 records the smallest
proved objective divided by $2\cdot10^8$ is
\[
 \frac{65955289561887369906031994339}
 {8192000000000000000000000000000}>\frac{805103}{10^8}.
\]
The exact positive margin is
\[
 \frac{1251801887369906031994339}
 {8192000000000000000000000000000}.
\]
These computations are finite computational evidence in the explicitly
stated integer and rational trust scope; they are not consequences of
floating-point solver status. The span-eight variable has no positive
cut in this certificate, so the joint old 33 pair terms already suffice.

\begin{proposition}[Uniform local improvement]\label{prop:nine}
Under the admitted continuous PC8 atom and capture semantics,
$F_9(g)\ge c=805103/10^8$ for every $g\in[4/5,\infty)^8$.
\end{proposition}
\begin{proof}
Outside the two-frame low-slack case, the first observation in this
section pays $c$. Inside it, completeness of the reflected outer cover
places the actual point in at least one of the 289 labelled
intersections. Its actual gaps and squares satisfy the associated
matrix and box. The exact dual bound pays $c$ there. This covers the
entire unbounded separated domain, including every closed boundary.
\end{proof}

\begin{lemma}[Direct sliding]\label{lem:sliding}
For any $m$ real points separated by at least $\theta$,
\begin{equation}\label{eq:sliding}
 \sum_{i\ne j}K(y_i-y_j)^2
 \ge c(m-8)-B\operatorname{span}(y).
\end{equation}
The span is zero when $m\le1$.
\end{lemma}
\begin{proof}
For $m\ge9$ sum Proposition~\ref{prop:nine} over the $m-8$
nine-point windows. Each pair of a fixed index span receives total
nonnegative weight at most two. Each adjacent gap receives total
fee at most $B$. The pair sum is at most the full ordered pair
sum, and the total pressure is at most $B\operatorname{span}(y)$.
For $m\le8$ the right side is nonpositive.
\end{proof}

\section{A band-limited majorant for the same window}

We use $\sinc t=\sin(\pi t)/(\pi t)$, with $\sinc0=1$, and put
\begin{equation}\label{eq:majorant}
 \theta=\frac45,\quad \gamma=\frac{61}{100},\qquad
 g(u)=\gamma\{\sinc(\theta(u-1/2))+
                    \sinc(\theta(u+1/2))\}^2.
\end{equation}

\begin{lemma}\label{lem:majorant}
The nonnegative function $g$ satisfies $g\ge f$ on $\R$,
$\widehat g(t)=0$ for $|t|\ge\theta$, and
\begin{equation}\label{eq:mass}
 \Lambda:=\int_\R g(u)\,du
 =\frac{2\gamma}{\theta}(1+\sinc\theta)
 \le\frac{61}{32}<2 .
\end{equation}
\end{lemma}
\begin{proof}
Each shifted sinc has a rectangular Fourier transform supported in
$[-\theta/2,\theta/2]$. The square in \eqref{eq:majorant}
therefore has Fourier support $[-\theta,\theta]$.
The convolution is continuous and zero at both boundary points,
so the conclusion includes equality $|t|=\theta$.
Parseval gives the displayed formula for $\Lambda$.
As $\sin(4\pi/5)=\sin(\pi/5)\le\pi/5$,
$\sinc(4/5)\le1/4$, proving its stated upper bound.

The remaining assertion $g\ge f$ is proved by the small exact
rational certificate in Appendix~\ref{app:certificate}.
Specifically it establishes
\[
 g(k/256)-f(k/256)>9/100 \quad(0\le k\le128),
 \qquad \sup_{u\in I}|(g-f)'(u)|<16.
\]
Both functions are even, and every point of $[0,1/2]$ is at
distance at most $1/512$ from one of these nodes. Hence on $I$
\[
 g(u)-f(u)>9/100-16/512=47/800>0.
\]
Outside $I$ the assertion follows from $f=0$ and $g\ge0$.
The derivative estimate is the essential continuous bridge;
the grid inequalities alone would not suffice.
\end{proof}

\begin{lemma}[Close-pair reward]\label{lem:close}
For $|t|\le4/5$, one has $K(t)>134/825$ and
$K(t)^2>1/40>c$.
\end{lemma}
\begin{proof}
For $0\le t\le4/5$, the functions $\cos(\sqrt2u)$ and
$\cos(2\pi tu)$ both decrease on $[0,1/2]$.
Chebyshev's integral covariance inequality, divided by $Z_0$,
gives for the unperturbed kernel $K_0(t)\ge\sinc t\ge\sinc(4/5)$.
The last monotonicity follows by differentiating $\sin x/x$ on
$[0,\pi]$. Using $\pi^2<10$ and
$\sin x\ge x-x^3/6$ at $x=\pi/5$ gives
$\sinc(4/5)>7/30$.
The perturbation of the normalized kernel has modulus at most
$S/Z_0<39/550$. Thus
$K(t)>7/30-39/550=134/825$.
Finally $(134/825)^2=17956/680625>1/40>c$.
Evenness covers negative $t$.
\end{proof}

\section{Lossless spectral assembly}

Choose even $\chi_\varepsilon\in C_c^\infty((-1/2,1/2))$,
$0\le\chi_\varepsilon\le1$, tending to one on the interior, and set
\[
 m_\varepsilon=\int_I f\chi_\varepsilon^2,\qquad
 f_\varepsilon=f\chi_\varepsilon^2/m_\varepsilon,\qquad
 K_\varepsilon=\widehat f_\varepsilon,\qquad
 d_\varepsilon=\|f_\varepsilon-f\|_1.
\]
Then $m_\varepsilon\to1$, $d_\varepsilon\to0$, and
$f_\varepsilon\to f$ in $L^1\cap L^2$.
The square root $\eta_\varepsilon=\sqrt{f_\varepsilon}$ is
smooth, even, compactly supported in the interior of $I$.
On the real axis, $|K_\varepsilon-K|\le d_\varepsilon$
and both kernels have modulus at most one.
Throughout the argument fix $\varepsilon$ sufficiently small that
$m_\varepsilon>61/64$.

For a finite set $Y$ of $n$ real points let
$U=(K_\varepsilon(y_i-y_j))_{i,j}$ be its positive Gram matrix,
with unit diagonal. Define
\[
 \phi_2(t)=t^2-(t-2)_+^2,\qquad
 J(U)=\tr\phi_2(U)-n .
\]
The scalar function $\phi_2$ is convex on $[0,\infty)$.

\begin{proposition}\label{prop:lossless}
For every finite $Y$,
\begin{equation}\label{eq:gain}
 J(U)\ge cn-B\,\operatorname{span}(Y)-8c-32d_\varepsilon n.
\end{equation}
\end{proposition}
\begin{proof}
Select a maximal collection of disjoint pairs with distance
strictly less than $\theta$. The remaining set $S$ is
$\theta$-separated. For any coefficient vector $(z_j)$ on $S$,
the pointwise majorant gives
\[
 \int f_\varepsilon(u)\left|\sum_{j\in S}z_je^{-2\pi iy_ju}\right|^2du
 \le\frac1{m_\varepsilon}
 \int g(u)\left|\sum_{j\in S}z_je^{-2\pi iy_ju}\right|^2du
 =\frac{\Lambda}{m_\varepsilon}\sum_{j\in S}|z_j|^2 .
\]
The cross terms vanish even at separation exactly $\theta$.
Thus $0\preceq U_{SS}\preceq(\Lambda/m_\varepsilon)I\prec2I$.
Normalization by $m_\varepsilon$ is necessary here.
Consequently $\phi_2(U_{SS})=U_{SS}^2$.

Write $m=|S|$. Apply each local certificate to $K$ and transfer
only its finitely many weighted pair terms to $K_\varepsilon$.
Each square changes by at most $2d_\varepsilon$; the total
weight in one nine-point window is at most $16$ by the span budgets.
The sliding argument, including the cases $m\le8$, yields
\[
 \tr\phi_2(U_{SS})-m
 \ge cm-B\,\operatorname{span}(S)-8c-32d_\varepsilon m .
\]
Each removed pair has a Gram matrix with eigenvalues
$1\pm|K_\varepsilon(t)|$ in $[0,2]$.
Its contribution to $J$ is
$2K_\varepsilon(t)^2\ge2c-4d_\varepsilon$, by
Lemma~\ref{lem:close}.

Diagonalize the principal block for $S$ and each two-point block.
Scalar Jensen in the resulting orthonormal basis proves
$\tr\phi_2(U)\ge\sum_{\text{blocks}}\tr\phi_2(U_{\text{block}})$.
This requires scalar convexity, not operator convexity.
Summing the pair rewards and the separated-set reward gives
\eqref{eq:gain}, since $\operatorname{span}(S)\le
\operatorname{span}(Y)$ and the pair error is at most
$2d_\varepsilon(n-m)\le32d_\varepsilon(n-m)$.
Empty or small separated sets are covered by the same single
boundary charge $8c$.
\end{proof}

\section{All-zero Hermitian counting and the energy constant}

\subsection{The unconditional analytic input}

For each fixed $\varepsilon$, the unconditional correlation formula
of Baluyot, Goldston, Suriajaya and Turnage-Butterbaugh
\cite[Lemma~5]{BGST}, with the correction described in
\cite[Section~3]{BGSTcorr}, gives
\begin{equation}\label{eq:unweighted}
 \sum_{0<\Im\rho,\Im\rho'\le T}
 K_\varepsilon\left(i(\rho-\rho')\frac{\log T}{2\pi}\right)^2
 =(C(f_\varepsilon)+o_\varepsilon(1))N(T),
\end{equation}
where
\begin{equation}\label{eq:energy}
 C(h)=\int_Ih(u)^2du+
           \int_I\int_I|u-v|h(u)h(v)\,du\,dv.
\end{equation}
The sum includes multiplicities and all off-line zeros.
For completeness, the weighted lemma applied to a fixed real even
test $q$ supported in $[-1,1]$ has weight
$w(\rho-\rho')=4/(4-(\rho-\rho')^2)$ and main term
$q(0)+2\int_0^1u q(u)\,du$.
Apply it separately to the two fixed smooth functions
$q=f_\varepsilon*f_\varepsilon$ and $q''$.
For $z=i(\rho-\rho')\log T/(2\pi)$, integration by parts gives
\[
 \left(\widehat q(z)-\frac{\widehat{q''}(z)}
                            {4(\log T)^2}\right)w(\rho-\rho')
 =K_\varepsilon(z)^2.
\]
The second fixed-test contribution is $O_\varepsilon(N/(\log T)^2)$;
the first gives \eqref{eq:unweighted}. The lemma allows signed tests,
so its use for $q''$ is legitimate.
This is the fixed-smoothing argument also used in \cite{Lamzouri}.
No estimate uniform in a $T$-dependent smoothing parameter is used.

\subsection{Retaining off-line zeros}

Put $L=\log T$ and $z_\rho=-i(\rho-1/2)L/(2\pi)$.
The multiset of these points is conjugation-invariant.
First in the complex Hilbert space $L^2(I;\mathbb C)$ set
$v_z(u)=\eta_\varepsilon(u)e^{-2\pi izu}$.
For each distinct nonreal conjugate pair write
$g_z=(v_z+v_{\bar z})/2$ and
$h_z=(v_z-v_{\bar z})/(2i)$.
The vectors $v_x$ for real $x$, together with $g_z$ and $h_z$,
belong to the real Hilbert space
$\mathcal H_{\R}=\{v:\overline{v(u)}=v(-u)\}$:
the conjugate reflection $J$ satisfies $Jv_z=v_{\bar z}$.
Define the finite self-adjoint operator on their real span by
\[
 A_T=\sum_{\text{distinct real }x}\mu_xv_x\otimes v_x+
       2\sum_{\text{distinct nonreal pairs}}\mu_z
                  (g_z\otimes g_z-h_z\otimes h_z).
\]
It acts on the finite span of these real-type vectors.
The identity $\|g_z\|^2-\|h_z\|^2=1$ and direct tensor expansion give
\begin{equation}\label{eq:operator}
 \tr A_T=N(T),\qquad
 \tr A_T^2=(C(f_\varepsilon)+o_\varepsilon(1))N(T).
\end{equation}
Indeed the Hilbert--Schmidt square expands to
$\sum_{z,w}K_\varepsilon(z-\bar w)^2$; conjugation reindexes this
as the sum in \eqref{eq:unweighted}. Individual complex terms
need not be positive.

Let $P$ be the sum of the unit features at simple real points,
$n=N_0^s(T)$, and $Q=A_T-P$. If $r$ is the number of repeated
real points and $k$ the number of nonreal conjugate pairs, then
\[
 n_+(Q)\le b:=r+k,\qquad N(T)\ge n+2b.
\]
Each remaining real point has multiplicity at least two; each
nonreal pair contributes at least two zeros and at most one
positive rank direction. No location-based deletion has occurred.

\begin{lemma}[Finite spectral counting]\label{lem:counting}
If $A=P+Q$, $P$ is the sum of $n$ unit features with Gram $U$,
$n_+(Q)\le b$, $\tr A=N$ and $N\ge n+2b$, then
\begin{equation}\label{eq:counting}
 \tr A^2\ge2(N-n)+\tr\phi_2(U)
           =2N-n+J(U).
\end{equation}
\end{lemma}
\begin{proof}
Enlarge the finite ambient space by zero directions if necessary.
If $\lambda_i(P)=p_i$ are its $n$ Gram eigenvalues, including
zeros, minmax gives $\lambda_{b+i}(A)\le p_i$.
For $t\le p$, $p\ge0$,
$4t-t^2\le4p-\phi_2(p)$; also $4t-t^2\le4$ for all real $t$.
After the first $b$ eigenvalues use the first inequality; any
remaining eigenvalues are nonpositive and contribute at most zero.
Since $\sum p_i=n$, this proves
$4N-\tr A^2\le4b+4n-\tr\phi_2(U)$.
Now $4b\le2(N-n)$ gives \eqref{eq:counting}.
Zeros in the $n$-dimensional Gram spectrum are included once;
the added ambient zero directions contribute nothing.
\end{proof}

\subsection{Energy and completion of the proof}

Define $(\mathcal Th)(u)=h(u)+\int_I|u-v|h(v)\,dv$.
For $f_0=\cos(\sqrt2u)/Z_0$, $(\mathcal Tf_0)''=0$ and evenness
makes $\mathcal Tf_0$ constant. Therefore a zero-mass perturbation
has zero cross energy with $f_0$. Orthogonality of the integer
cosines, and
$\mathcal T\cos(2\pi ju)=(1-1/(2\pi^2j^2))\cos(2\pi ju)
+\text{constant}$, give
\begin{equation}\label{eq:amenergy}
 C(f)=\frac12+\beta\cot\beta+
 \frac{\displaystyle\sum_{j=1}^{12}c_j^2
           (1/2-1/(4\pi^2j^2))}{2\sin^2\beta}.
\end{equation}
The exact rational computation in Appendix~\ref{app:certificate}
establishes
\begin{equation}\label{eq:gainconstant}
 2-C(f)>\frac{67216841}{10^8}.
\end{equation}
In particular the energy cost of the perturbation is paid;
the unperturbed energy constant is not substituted.

\begin{proof}[Proof of Theorem~\ref{thm:main}]
The simple real positions lie in $(0,X_T]$, with
$X_T=T\log T/(2\pi)$ and $X_T/N(T)\to1$ by
Riemann--von Mangoldt. Apply Proposition~\ref{prop:lossless}
to their Gram and then Lemma~\ref{lem:counting}:
\[
 (1-c+32d_\varepsilon)n
 \ge2N(T)-\tr A_T^2-BX_T-8c.
\]
Fix sufficiently small $\varepsilon$, let $T\to\infty$ using
\eqref{eq:operator}, and only then let $\varepsilon\to0$.
The $L^1\cap L^2$ convergence implies
$C(f_\varepsilon)\to C(f)$ in \eqref{eq:energy}.
Consequently
\[
 \liminf_T\frac{n}{N(T)}
 \ge\frac{2-C(f)-B}{1-c}
 >\frac{67216841-404350}{10^8-805103}
 =\frac{66812491}{99194897}.
\]
The weak inequality stated in the theorem is therefore valid.
The boundary charge $8c$ disappears only after division by $N(T)$.
\end{proof}

\section{Verification scope and further improvement}\label{sec:trust}

Input~\ref{in:pc8} is pinned to the public
\texttt{dani-bound-2026} source \cite{PC8} at commit
\[
\texttt{d272437e7ebeb17b87c8c3d7cece592aa23785ab}.
\]
Its raw \texttt{Solution.lean} is 1,675,641 bytes with SHA-256
\[
\texttt{\footnotesize 012c6ac5f9282158a686500dc0c967bf0a9b00e1a6192734d8f237bb23dd2d5f}.
\]
The repository admission combines a continuous interval-soundness
audit with all 240 exact integer compiled evaluations, seven complete
finite structural checks, and Lean kernel-checked generic structural
and minorant bridges. It includes the value and derivative tables,
point lists, convexity chains, noncompact gap coverage and all 32
root trees. The evidence is recorded in repository note~483
and its linked audits \cite{Repo}.
The compiled evaluations are not 240 kernel-reduced proofs, and
the full upstream zeta formalization was not rebuilt.
The present result has that explicit finite-computation trust scope.

The new majorant and energy checker uses only Python's standard-library
arbitrary-precision integers and rational arithmetic, with assertions
enabled. It does not run the much larger native search supporting
the numerical headline in \cite{Knaus}. In particular that search's
\texttt{native\_decide} computation axiom is not an input to this paper.
Neither this checker nor the earlier symbolic bridges constitute
a complete Lean verification of Theorem~\ref{thm:main}.
The mathematical transfer, finite computation and formalized
sublemmas are separate evidence.

The new finite proof is recorded by
{\small\path{scripts/am_nine_point_epigraph_certificate.py}} and
{\small\path{output/am-nine-point-epigraph-certificate.json}}.
The latter includes all 70 captured cells with their 26 atom records,
16 decoded Farkas integers, and all 289 sparse rational duals.
The verifier reconstructs the original minorant guard before consuming
the new cuts. Its \texttt{--replay} mode rebuilds the isolated capture
from the fixed raw source and the two admitted generic bridges,
checks all 41 old outer/root results, and compares every captured record.
Its \texttt{--check} mode consumes saved records only. Both modes were
actually run with assertions enabled; optimized Python is refused.
The full earlier value and derivative admission is reused, rather
than falsely described as repeated by the new 41 evaluations.
The new certificate and analytic transfer were independently reviewed
inside the project. Internal AI-assisted review is distinct from
external peer review or complete kernel formalization.

Further improvement requires a stronger continuous local reward,
a better window, or an actually paid additional spectral gain.
The fixed-window method ceiling is not a bound on the true zero
proportion. The present result does not improve the project's
separate zero-free-region bounds or pay its remaining original
four-prime arithmetic covariance.

\appendix
\section{Small exact rational certificate}\label{app:certificate}

The accompanying standard-library program is
{\small\path{scripts/am_lossless_majorant_certificate.py}}, with report
{\small\path{output/am-lossless-majorant-certificate.json}}.
From the repository root run
\[
 \texttt{python -B scripts/am\_lossless\_majorant\_certificate.py --check}.
\]
The checker refuses optimized Python execution.
The following specifies the full continuous bridge independently
of floating-point output.

For a nonnegative rational interval $q\subset[0,10)$, it encloses
$\cos\sqrt q$ and $\sin\sqrt q/\sqrt q$ by interval evaluation of
\[
 \sum_{k=0}^{12}\frac{(-1)^kq^k}{(2k)!},
 \qquad
 \sum_{k=0}^{12}\frac{(-1)^kq^k}{(2k+1)!},
\]
with symmetric remainder bounds
$q_{\max}^{13}/26!$ and $q_{\max}^{13}/27!$, respectively.
The alternating tails decrease from the first omitted term.
For the integer cosine modes the rational phase $2ju$ is
reduced modulo two to $[-1,1]$ before multiplication by $\pi$.
No division by a small trigonometric argument is required.
The normalization $Z_0=\sin\beta/\beta$ has the exact argument
squared $q=1/2$.

Machin's identity
$\pi=16\arctan(1/5)-4\arctan(1/239)$, with the first 24 terms
of each alternating arctangent series and its next-term remainder,
certifies
\[
 3141592653/10^9<\pi<3141592654/10^9.
\]
These enclosures, at all 129 nodes $u=k/256$, certify the strict
lower margin $g(u)-f(u)>9/100$.
On all of $I$, $|\sinc'|\le\pi/2$ follows from its integral
representation. Thus
\[
 |(g-f)'|\le4\gamma\pi\theta+
       \frac{\sqrt2+2\pi\sum_jj|c_j|}{Z_0}
 <4\gamma\frac{22}{7}\theta+
       \frac{3/2+(44/7)\sum_jj|c_j|}{9/10}<16.
\]
This proves the global margin $47/800$ stated above.

For \eqref{eq:gainconstant}, use the same Taylor intervals
to bound $\cos\beta/Z_0$ from above, $Z_0^2=2\sin^2\beta$
from below, and $\pi$ from above in the positive correction
in \eqref{eq:amenergy}. All resulting comparisons are exact rational
inequalities. A shorter independent certificate uses
\[
\begin{aligned}
 u&=\frac{3329448031}{4379443200}>\cos\beta,&
 z&=\frac{153778383027779}{167382319104000}<Z_0,\\
 d&=\frac{180493}{213840}<2\sin^2\beta,&
 p&=\frac{5277328977275528}{1679825970703125}>\pi.
\end{aligned}
\]
These arise, respectively, from cosine terms through index 6,
sinc terms through index 7, cosine at $\sqrt2$ through index 6,
and Machin arctangent upper/lower sums through indices 4 and 1.
Direct rational arithmetic gives
\[
 \frac32-\frac uz-
 \frac{\sum_jc_j^2(1/2-1/(4p^2j^2))}{d}
 >\frac{67216841}{10^8}.
\]
The positive margin is greater than $5.6\cdot10^{-10}$.
The old checker also verifies the old span capacities, close-pair
constants and majorant mass bound. The new rational ratio is checked
by the accompanying nine-point verifier; the joint span budgets follow
directly from \eqref{eq:capacity} and the definition of $a'$.
These small checks do not substitute for
the full seven-dimensional certificate in Input~\ref{in:pc8}.

\begin{thebibliography}{9}
\bibitem{BGST}
S. A. C. Baluyot, D. A. Goldston, A. I. Suriajaya and
C. L. Turnage-Butterbaugh,
\emph{An unconditional Montgomery theorem for pair correlation of
zeros of the Riemann zeta function},
Acta Arith. 214 (2024), 357--376;
\href{https://arxiv.org/abs/2306.04799}{arXiv:2306.04799}.

\bibitem{BGSTcorr}
S. A. C. Baluyot, D. A. Goldston, A. I. Suriajaya and
C. L. Turnage-Butterbaugh,
\emph{Pair correlation of zeros of the Riemann zeta function I:
proportions of simple zeros and critical zeros},
\href{https://arxiv.org/abs/2501.14545v3}{arXiv:2501.14545v3},
Section~3 and its correction of the earlier Theorem~1.

\bibitem{Lamzouri}
Y. Lamzouri,
\emph{A new proof that more than $2/3$ of the zeros of the Riemann
zeta function are simple and on the critical line},
\href{https://arxiv.org/abs/2609.02882v1}{arXiv:2609.02882v1},
Section~3, fixed-smoothing removal of the correlation weight.

\bibitem{PC8}
\emph{AM eight-point PC8CL certificate}, public
\texttt{dani-bound-2026} submission, fixed commit
\texttt{d272437e7ebeb17b87c8c3d7cece592aa23785ab},
\href{https://github.com/josusanmartin/riemann/tree/d272437e7ebeb17b87c8c3d7cece592aa23785ab/submissions/dani-bound-2026/proof}{public proof source}.
The cosine-perturbed window and local weights are prior inputs.

\bibitem{Knaus}
K. M. Knausg{\aa}rd,
\emph{More than 83.9\% of the zeros of the Riemann zeta function
are distinct and more than 67.35\% are simple and on the critical line},
\href{https://arxiv.org/abs/2610.08965v1}{arXiv:2610.08965v1},
October 6, 2026, especially Section~5.

\bibitem{Lossless}
RH--Weil research project,
\emph{Lossless Assembly of an Eight-Point Certificate for Simple
Critical Zeros of the Riemann Zeta Function}, October 8, 2026,
\href{https://github.com/lixiang90/RH-Weil/blob/0f7e6c21769b7a848fc9ce4cbf498eb91184874d/papers/lossless-eight-point-simple-critical-paper.tex}{fixed manuscript and proof}.

\bibitem{Repo}
RH--Weil research project,
\href{https://github.com/lixiang90/RH-Weil}{public repository},
note~483, the PC8 finite replay and independent continuous audits;
notes~498--501 and the joint nine-point certificate.
\end{thebibliography}
\end{document}
